\section{신뢰 경계와 재현 방법}
\label{sec:trust-reproduction}

\subsection{무엇을 신뢰하는가}

이 증명 개발의 신뢰 기반(trusted computing base, TCB)은 다음을 포함한다.

\begin{itemize}
  \item EasyCrypt 커널/형 검사기(kernel/type checker)와 가져온 표준 라이브러리;
  \item Why3 및 호출된 SMT 증명기(prover);
  \item Jasmin 구문분석기/컴파일러(parser/compiler)와 \identifier{jasmin2ec}
        번역기;
  \item 해시(hash)로 고정한 Jasmin 소스, 생성 추출물, 재사용 EasyCrypt 산출물;
  \item 소스 해시(source hash) 및 LaTeX 빌드를 수행하는 호스트 도구.
\end{itemize}

\identifier{jasmin2ec}가 소스를 받아들이고 생성된 EasyCrypt가 컴파일된다는
사실은 번역기의 의미보존 증명이 아니다. 별도의 검증된 번역기 정리가 없는 한
번역기 자체는 TCB에 남는다. 마찬가지로 메모리 안전성과 SCT/상수시간
(constant-time)은
EUF-CMA에서 따라오지 않으며 별도의 트랙이다.

\subsection{증명 검증 파이프라인}

저장소 루트에서 다음을 실행한다.

\begin{verbatim}
./haetae-topdown-easycrypt/scripts/verify-all.sh
\end{verbatim}

이 스크립트는 다음을 순서대로 검사한다.

\begin{enumerate}
  \item 고정 소스(pinned source)와 읽기 전용 증명 트리(read-only proof tree)의
        SHA-256 변동(drift);
  \item 패커(packer), 서명 코덱, HBZ 준비/적용/전체 래퍼, API 키 복사 보조절차,
        원시 API 호출자, 원시 \(\mu\), 전사 집중 추출(focused extraction)의 재생성과 hash
        일치, Week~7 pack/unpack과 Week~8 focused HBZ procedure body identity 및
        Week~9 실제 rANS 결합 모듈(actual rANS harness)이 재사용하는 생성 대상
        동일성(generated-target identity);
  \item 실제 512단어 HBZ 기호 인증서(symbol certificate)의 결정적 재생성,
        생성기/출력(generator/output) 해시 일치와 생성된 EasyCrypt 인증서 자체의
        커널(kernel) 재검사;
  \item Week~9 rANS 정리 표면에 더해 Week~10 배열/목록, 단어 단계, 내부 진행,
        실제 내부 반복문, 직렬화와 성공 크기 정리 표면 및 Week~11 꼬리 인식
        불변식, 생성 단어 단계, 최종 저장과 실제 인코더 추적 폐쇄, Week~12
        디코더 커서/표/재생 불변식과 실제 추적 정제, Week~13 복사 추적/점별
        입력/W64 크기 어댑터와 실제 핵심 역함수의 과대주장 검사;
  \item 현재 매니페스트된 67개 작성 EasyCrypt 대상의 \identifier{-no-eco} 새
        컴파일;
  \item \identifier{admit}, \identifier{sorry}, 작성 \identifier{axiom},
        디버그/임시(debug/temporary) 선언 부재와 매니페스트(manifest) 완전성;
  \item 선택된 상류 기준선(upstream baseline)의 새 컴파일(fresh compile);
  \item 기존 주차별(sprint) LaTeX 연구 노트의 참조/인용 오류 없는 빌드.
\end{enumerate}

추출 해시(hash)는 생성물 변동(drift)을 탐지하지만 번역 의미의 정당성을
대신하지 않는다. 또한 실행 요약은 매 실행 때 덮어쓰므로, 발표용 스냅샷에는
마지막 \texttt{RESULT PASS} 줄까지 포함된 완전한 로그를 보관해야 한다.
\sourcepath{WEEK13_REPORT.md}가 기록한 Week~13 완료 실행의 종료 표식은
\texttt{RESULT PASS authored-targets=67}\newline
\texttt{cache=-no-eco}다.

\begin{boundary}[67대상 완료 스냅샷]
2026-08-09의 Week~13 완료 실행은 Week~12의 65대상 기준선에
\sourcepath{Mode2RansCoreCompositionBridge.ec}와
\sourcepath{Mode2RansCoreActualInverse.ec}를 더한 67개를 모두 컴파일한 뒤
기준선, 주차별 LaTeX와 읽기 전용 소스 무변경 검사를 통과했다. 독립 검증 기록
\sourcepath{logs/week13-independent-verifier.md}도 최종 정리의 사전조건에
인코더 성공, 복사 추적, 디코더 성공이나 출력 등식이 숨어 있지 않음을
확인한다.
실행 중 생성되는 \sourcepath{logs/verify-all-summary.txt}는 마지막 줄이
\texttt{RESULT PASS}일 때만 완료 실행의 증거다. 그 줄이 없는 중간 또는
중단 실행은 보존된 완료 기준선을 갱신하지 않는다.
\end{boundary}

\subsection{이 문서 빌드}

이 안내서는 기존 sprint 노트와 독립된 XeLaTeX 문서다. 저장소 루트에서

\begin{verbatim}
sh haetae-topdown-easycrypt/theory-guide/build.sh
\end{verbatim}

을 실행하면 \sourcepath{haetae-topdown-easycrypt/theory-guide/main.pdf}가
생성된다. 스크립트는 중단 오류와 정의되지 않은 참조/인용(undefined
reference/citation)을 실패로 처리한다.

문서 작성 시 관찰된 도구 경계는 Jasmin/\identifier{jasmin2ec} 2026.03.0,
EasyCrypt OPAM switch \identifier{easycrypt-5.2}, Why3 1.8.0이며, 구성된
prover에는 Alt-Ergo 2.6.0, CVC4 1.8.0, CVC5 1.2.1, Z3 4.12.6이 포함된다.
정확한 재현 결과는 버전 문자열보다 실행 로그를 우선한다.

\begin{takeaway}{완료 판정 기준}
이 안내서가 컴파일되는 것과 EasyCrypt 정리가 컴파일되는 것은 서로 다른
증거다. 문서 배포 전에는 \identifier{build.sh}와 전체
\identifier{verify-all.sh}를 모두 성공시켜야 한다. 현재 증명 기준선은
67대상 전체 통과이며, 이 안내서의 별도 빌드도 함께 확인한다.
\end{takeaway}
